% !PRINTED 287
\weberpage{287}
Wenn nun die positive Fortschrittsrichtung auf $\varphi$ an einem
der Durchschnittspunkte der Curven $\varphi$, $\psi$ in das Innere von $\psi$
hineinführt, so ist $d\psi$, also auch die Functionaldeterminante
$[\varphi, \psi]$ negativ, und im entgegengesetzten Falle positiv (Fig.~7).
Wir nennen einen Schnittpunkt von $\varphi$ und $\psi$ einen
\gesperrt{Austrittspunkt} $A(\varphi, \psi)$ oder einen Eintrittspunkt
$E(\varphi,\psi)$, je nachdem

\weberfigureimage{weber-b1-fig06-07.png}{0.94}

die positive Fortschrittsrichtung von $\varphi$ aus dem Inneren von $\psi$
ins Aeussere oder vom Aeusseren ins Innere führt, und haben also den Satz:
\weq{4}{
\begin{array}{lll}
\text{an einem } A(\varphi, \psi) & \text{ist} & [\varphi, \psi] > 0 \\
\text{,,}\ \ \text{,,} & E(\varphi, \psi) & \text{,,}\ \ [\varphi, \psi] < 0.
\end{array}
}
Die Anzahl der $A(\varphi, \psi)$ ist ebenso gross, wie die der $E(\varphi, \psi)$;
und wenn wir $\varphi$ mit $\psi$ vertauschen, so gehen die $A(\varphi, \psi)$, $E(\varphi, \psi)$
in $E(\psi, \varphi)$ und $A(\psi, \varphi)$ über.

Die zwei Functionen $\varphi$, $\psi$ bestimmen in eindeutiger Weise
einen Flächenraum, der dadurch charakterisirt ist, dass in ihm
das Product $\varphi\psi$ negativ ist. Wir wollen ihn den \gesperrt{Binnenraum}
$(\varphi, \psi)$ oder $(\psi, \varphi)$ nennen. In unserer Fig.~7 ist es die schraffirte
Fläche.

\begin{center}
\S.\,94.

\gesperrt{Charakteristik eines Systems von drei Functionen.}
\end{center}

Wir nehmen nun zu den beiden Functionen $\varphi$, $\psi$ eine dritte
$f(x, y)$ hinzu. Die durch die Gleichung $f = 0$ dargestellte
Curve, oder die Curve $f$ soll gleichfalls im Endlichen geschlossen
%CONT>>
